\documentclass[11pt]{article}
\usepackage[margin=1in]{geometry}
\usepackage{amsmath,amssymb,hyperref}
\title{Atom 2.1: Electron Closure and the Fine-Structure Charge--Action Interface}
\author{WeBeGood [author list and affiliation to be confirmed]}
\date{September 11, 2026}
\begin{document}\maketitle
\begin{abstract}
A conditional reduction, not a numerical prediction of the fine-structure
constant. The full canonical derivation and limitations are in paper.md and
nodes N017/N018. This compact LaTeX counterpart records the central equations.
\end{abstract}
\section{Closure family}
For a candidate torus trefoil around the seed-line x axis, let
\begin{align}
x&=a\sin3t,&y&=(R+a\cos3t)\cos2t,&z&=(R+a\cos3t)\sin2t,\\
L_e&=R F(\rho),&F(\rho)&=\int_0^{2\pi}\sqrt{4(1+\rho\cos3t)^2+9\rho^2}\,dt.
\end{align}
Here $\rho=a/R\in(0,1)$. With the pitch convention $L_\nu=2\pi/|\kappa|$
and $R=\beta L_\nu$, $L_e=\beta F(\rho)L_\nu$. An additional scalar phase law
$|\kappa|L_e=2\pi n$ implies $\beta=n/F(\rho)$, but selects neither shape nor
stable field dynamics. Curl eigenvalue and phase transport are not identical
without an extra assumption.
\section{Charge and action}
For smooth source-free fields, $\nabla\cdot\mathbf E=0$ implies
$Q=\epsilon_0\oint\mathbf E\cdot d\mathbf A=0$. Extensions require explicit
sources, medium response or boundary/topology accounting.
For a finite-energy candidate set $\boldsymbol\xi=\mathbf r/L$,
$\mathbf E=E_*\mathbf f$, $\mathbf B=E_*\mathbf g/c$ and
\begin{align}
C_Q&=\frac{1}{4\pi}\lim_{s\to\infty}\oint_{|\xi|=s}\mathbf f\cdot d\mathbf A_\xi,&
C_U&=\frac12\int (|\mathbf f|^2+|\mathbf g|^2)d^3\xi,\\
Q&=4\pi\epsilon_0 E_*L^2 C_Q,&U_{\rm EM}&=\epsilon_0 E_*^2L^3 C_U.
\end{align}
If $|Q|=e$ and $\eta_{\rm EM}=U_{\rm EM}L/(\hbar c)$, then
\begin{equation}\boxed{\alpha_{\rm fs}=4\pi\eta_{\rm EM}\frac{C_Q^2}{C_U}.}\end{equation}
If total energy includes a medium/core contribution, use
$\eta_{\rm EM}=f_{\rm EM}\eta_{\rm total}$. A future independently justified
$U_{\rm total}/\omega=\hbar$ would imply $\eta_{\rm total}=\omega L/c$;
no such normalization is established here.
\section{Diagnostic and nonuniqueness}
The supplied-source profile $\mathbf f=\boldsymbol\xi/(1+|\xi|^2)^{3/2}$,
$\mathbf g=0$ has $C_Q=1$, $C_U=3\pi^2/8$ and
$\alpha_{\rm fs}=32\eta_{\rm EM}/(3\pi)$. It is not an electron solution.
Amplitude freedom changes the coupling while preserving shape.
For nonzero baseline coefficients,
\begin{equation}
\frac{\delta\alpha}{\alpha}=\frac{\delta\eta_{\rm EM}}{\eta_{\rm EM}}
+2\frac{\delta C_Q}{C_Q}-\frac{\delta C_U}{C_U}+O(\delta^2).
\end{equation}
This is sensitivity bookkeeping, not derived correction coefficients.
\section{Units and completion status}
In SI, $\alpha=e^2/(4\pi\epsilon_0\hbar c)$; in Gaussian units,
$\alpha=e_G^2/(\hbar c)$; in Heaviside--Lorentz units with $\hbar=c=1$,
$\alpha=e_{HL}^2/(4\pi)$. The CODATA 2022 benchmark is
$\alpha^{-1}=137.035999177(21)$\cite{nist}; it is never a normalization input.
Charge-producing dynamics, stable closure, action selection, medium energy and
controlled low-energy corrections remain open. No numerical alpha is derived.
\section{Reproducibility and declarations}
Run the N017/N018 Python diagnostic programs and the repository tests. Full
assumptions, falsifiers, author/contribution placeholders, funding/conflict
placeholders and version provenance are in paper.md. WeBeGood supplied project
direction; ChatGPT/Codex assisted drafting and computation. No experimental data
or human/animal research were produced. Human approval is pending publication.
\bibliographystyle{plain}\bibliography{refs}
\end{document}
